\documentclass[CJK]{beamer}
\usepackage{CJKutf8}
\usepackage{beamerthemesplit}
\usetheme{Malmoe}
\useoutertheme[footline=authortitle]{miniframes}
\usepackage{amsmath}
\usepackage{amssymb}
\usepackage{graphicx}
\usepackage{color}
\usepackage{slashed}
\usepackage{simplewick}
\graphicspath{{../figures/}}
\def\be{\begin{equation}}
\def\ee{\nonumber\end{equation}}
\def\bea{\begin{eqnarray}}
\def\eea{\nonumber\end{eqnarray}}
\def\ii{{\dot{\imath}}}
\def\bch{\begin{CJK}{UTF8}{gbsn}}
\def\ech{\end{CJK}}
\def\bex{\begin{minipage}{0.3\textwidth}\includegraphics[width=1in]{jugelizi.png}\end{minipage}\begin{minipage}{0.6\textwidth}}
\def\eex{\end{minipage}}
\def\chtitle#1{\frametitle{\bch#1\ech}}
\def\skipline{{\vskip0.1in}}
\def\skiplines{{\vskip0.2in}}
\def\lagr{{\mathcal{L}}}
\def\hamil{{\mathcal{H}}}
\def\vecv{{\mathbf{v}}}
\def\vecx{{\mathbf{x}}}
\def\veck{{\mathbf{k}}}
\def\vecp{{\mathbf{p}}}
\def\vecn{{\mathbf{n}}}
\def\vecA{{\mathbf{A}}}
\def\vecP{{\mathbf{P}}}
\def\vecsigma{{\mathbf{\sigma}}}
\def\hatJn{{\hat{J_\vecn}}}
\def\hatJx{{\hat{J_x}}}
\def\hatJy{{\hat{J_y}}}
\def\hatJz{{\hat{J_z}}}
\def\hatj#1{\hat{J_{#1}}}
\def\hatphi{{\hat{\phi}}}
\def\hatq{{\hat{q}}}
\def\hatpi{{\hat{\pi}}}
\def\vel{\upsilon}
\def\Dint{{\mathcal{D}}}
\def\adag{{\hat{a}^\dagger}}
\def\bdag{{\hat{b}^\dagger}}
\def\cdag{{\hat{c}^\dagger}}
\def\ddag{{\hat{d}^\dagger}}
\def\hata{{\hat{a}}}
\def\hatb{{\hat{b}}}
\def\hatc{{\hat{c}}}
\def\hatd{{\hat{d}}}
\def\hatN{{\hat{N}}}
\def\hatH{{\hat{H}}}
\def\hatp{{\hat{p}}}
\def\Fup{{F^{\mu\nu}}}
\def\Fdown{{F_{\mu\nu}}}
\def\newl{\nonumber \\}
\def\SIkm{\mathrm{km}}
\def\SIyr{\mathrm{yr}}
\def\SIGyr{\mathrm{Gyr}}
\def\SIeV{\mathrm{eV}}
\def\SIGeV{\mathrm{GeV}}
\def\SIm{\mathrm{m}}
\def\SIcm{\mathrm{cm}}
\def\SIJ{\mathrm{J}}
\def\SIs{\mathrm{s}}
\def\SIkg{\mathrm{kg}}
\def\SIg{\mathrm{g}}
\def\vece{\mathrm{e}}
\def\bmat#1{\left(\begin{array}{#1}}
\def\emat{\end{array}\right)}
\def\bcase#1{\left\{\begin{array}{#1}}
\def\ecase{\end{array}\right.}
\def\calM{{\mathcal{M}}}
\def\calT{{\mathcal{T}}}
\def\calR{{\mathcal{R}}}
\def\barpsi{\bar{\psi}}
\def\baru{\bar{u}}
\def\barv{\bar{\upsilon}}
\def\bmini#1{\begin{minipage}{#1\textwidth}}
\def\emini{\end{minipage}}
\def\qeq{\stackrel{?}{=}}
\def\torder#1{\mathcal{T}\left(#1\right)}
\def\rorder#1{\mathcal{R}\left(#1\right)}


\title{Quantum Field Theory I \\ Homework 1}
  \author{}
  \date{}


\begin{document}

\begin{frame}
 
\begin{center}
\begin{Large}
\bch
量子场论 I 

{\vskip 0.3in}

第一次课后作业 （共八次，每次2.5分）

交作业时间: 9月12日，星期二，9：50am

\ech
\end{Large}
\end{center}

\vskip 0.2in

\bch
课件下载
\ech
https://github.com/zqhuang/SYSU\_QFTI

\end{frame}

\begin{frame}
\chtitle{第1题(0.5分）}
\bch
若算得的截面为$\sigma = 10^{-3}/m_W^2$, $m_W \approx 80 GeV$是$W^{\pm}$的粒子质量，试换算出以$\mathrm{cm}^2$为单位的截面值。若算得的寿命的$\tau = 100/m_w$，试问等于多少秒？
\ech
\end{frame}



\begin{frame}
\chtitle{第2题(0.5分）}
\bch
证明任意四维时空坐标系$(x^0, x^1, x^2, x^3)$的积分元$\sqrt{-g}d^4x$是一个标量。其中$g$是度规矩阵$g_{\mu\nu}$的行列式的简写，$d^4x$是积分元$dx^1dx^2dx^3dx^4$的简写。
\ech
\end{frame}


\begin{frame}
\chtitle{第3题(0.5分）}
\bch
考虑一维空间x和一维时间t构成的时空里的标量场$\phi(x,t)$，若其作用量为
$$ S = \int dx dt \,   \left[ \frac{1}{2}\left(\frac{\partial \phi}{\partial t}\right)^2 -  \frac{1}{2}\left(\frac{\partial \phi}{\partial x}\right)^2 - V(\phi)\right]\, ,$$
其中$V(\phi)$为给定的势能函数。试用求作用量稳定点的方法推导$\phi(x, t)$的运动方程并将结果与Euler-Lagrange方程做比较。
\ech
\end{frame}

\begin{frame}
\chtitle{第4题(0.5分）}
\bch
谐振子
$$ S = \int_{-\infty}^\infty \frac{m}{2}\left[(\frac{d\phi}{dt})^2 - \omega^2\phi^2\right] dt$$
把时间维特殊化以后就只有一个自由度$\phi$。写出$\phi$对应的正则动量$\pi$，系统的Hamilton量，以及Hamilton方程。试把$\pi$从两个Hamilton方程中消去，得到的结果和Euler-Lagrange方程一致吗？
\ech
\end{frame}


\begin{frame}
\chtitle{第5题(0.5分）}
\bch
作用量
$$S = \int d^4x\ \lagr(\phi, \partial_\mu\phi)\,.$$
在坐标平移下显然具有不变性。由此用Noether定理推导场$\phi$的能量动量守恒方程。

\skipline

提示：考虑四种无穷小变化$x^\mu \rightarrow x^\mu + \epsilon_r \delta_r^\mu$ ($r = 0, 1, 2, 3$)带来的$\phi$的变化和$\mathcal{L}$的变化。
\ech
\end{frame}


\end{document}
